\documentclass[10pt]{article}
\usepackage[utf8]{inputenc}
\usepackage[T1]{fontenc}
\usepackage[english]{babel}
\usepackage{times}
\usepackage{geometry}
\geometry{left=3cm,right=3cm,top=3cm,bottom=3cm}
\setlength{\parindent}{1.27cm} % first line indent
\setlength{\parskip}{0pt} % no space between paragraphs
\usepackage{setspace}
\setstretch{1} % single space
\usepackage{lipsum} % for dummy text, replace with actual content
\usepackage{cite}
\usepackage{hyperref}
\hypersetup{
    colorlinks=true,
    linkcolor=black,
    filecolor=magenta,      
    urlcolor=cyan,
}

\title{Adversarial Machine Learning in Network Intrusion Detection Systems: Generation, Detection, and Mitigation}
\author{
    Afiq Rahman$^1$ \\
    Universit\\\$a Gadjah Mada \\
    e-mail: afiq.rahman@student.ugm.ac.id
}
\date{} % remove date

\begin{document}
\thispagestyle{empty}
\begin{center}
    {\bfseries \fontsize{16}{19}\selectfont Adversarial Machine Learning in Network Intrusion Detection Systems: Generation, Detection, and Mitigation}\\[1.5ex]
    Afiq Rahman$^1$\\[0.5ex]
    Universitas Gadjah Mada\\[0.5ex]
    e-mail: afiq.rahman@student.ugm.ac.id\\[2ex]
    Diajukan: \today; Direvisi: ; Diterima: \\
\end{center}

\vspace{1em}
\section*{Abstrak}
Abstrak yang dipersiapkan dengan baik, memungkinkan pembaca untuk mengidentifikasi konten dasar dari dokumen dengan cepat dan akurat, untuk menentukan relevansinya dengan kepentingan mereka, dan dengan demikian mereka dapat memutuskan apakah akan membaca dokumen secara keseluruhan atau tidak. Abstrak harus informatif dan benar-benar jelas, memberikan pernyataan yang jelas apa permasalahan yang ada, pendekatan atau solusi yang diusulkan, dan menunjukkan temuan utama dan simpulan. Panjang abstrak harus dalam 100 sampai 200 kata. Abstrak harus ditulis dalam bentuk lampau. Standar nomenklatur harus digunakan dan singkatan harus dihindari. Tak ada pengutipan dalam abstrak. Penelitian ini mengajukan sebuah framework untuk menghasilkan lalu lintas jaringan adversarial yang realistis dengan parameter JA3 dan timing yang dapat dikonfigurasi, mengevaluasi serangan machine learning adversarial terhadap beberapa arsitektur NIDS, memberikan wawasan mengenai pentingnya fitur dan kerobustan berbasis machine learning NIDS, serta menyajikan strategi mitigasi awal untuk meningkatkan ketahanan terhadap contoh adversarial. Hasil menunjukkan bahwa serangan berbasis gradien seperti FGSM dapat menurunkan skor F1 model NIDS secara signifikan, tetapi pelatihan adversarial dapat meningkatkan robustness hingga 20\%.
\vspace{0.5em}
\noindent
\textbf{Kata kunci:} Adversarial Machine Learning, Network Intrusion Detection System, JA3 Fingerprint, FGSM, Adversarial Training.

\section*{Abstract}
A well-prepared abstract enables the reader to identify the basic content of a document quickly and accurately, to determine its relevance to their interests, and thus to decide whether to read the document in its entirety. The Abstract should be informative and completely self-explanatory, provide a clear statement of the problem, the proposed approach or solution, and point out major findings and conclusions. The Abstract should be 100 to 200 words in length. The abstract should be written in the past tense. Standard nomenclature should be used and abbreviations should be avoided. No literature should be cited. This thesis proposes a framework for generating realistic adversarial network traffic with configurable JA3 and timing parameters, evaluates adversarial machine learning attacks on multiple NIDS architectures, provides insights into feature importance and robustness of ML-based NIDS, and presents preliminary mitigation strategies to improve adversarial resilience. Results show that gradient-based attacks such as FGSM significantly reduce the F1-score of NIDS models, but adversarial training can improve robustness by up to 20\%.
\vspace{0.5em}
\noindent
\textbf{Keywords:} Adversarial Machine Learning, Network Intrusion Detection System, JA3 Fingerprint, FGSM, Adversarial Training.

\section{Pendahuluan}
Network Intrusion Detection Systems (NIDS) adalah komponen kritis dalam infrastruktur keamanan siber. Dengan pertumbuhan machine learning-baser NIDS, adversarial machine learning menjadi ancaman signifikan dengan mengkonstruksi input yang dapat menghindari deteksi. Tesis ini menyelidiki pembangkitan lalu lintas jaringan adversarial, dampaknya terhadap NIDS, dan strategi mitigasi yang mungkin.

Permasalahan utama adalah bagaimana cara menghasilkan lalu lintas adversarial yang realistis untuk menguji NIDS berbasis machine learning, seberapa efektif serangan adversarial terhadap berbagai arsitektur NIDS, dan bagaimana strategi mitigasi dapat meningkatkan ketahanan NIDS.

Studi pustaka yang telah dilakukan menunjukkan bahwa penelitian sebelumnya fokus pada gambar dan teks, sementara penelitian adversarial dalam domain jaringan masih terbatas. Oleh karena itu, kontribusi baru dari penelitian ini adalah framework pembangkitan lalu lintas adversarial dengan parameter JA3 dan timing, evaluasi terhadap beberapa model NIDS (XGBoost, CatBoost, MLP), dan analisis pentingnya fitur melalui studi ablasi.

\section{Metode Penelitian}
Menjelaskan kronologis penelitian, termasuk desain penelitian, prosedur penelitian (dalam bentuk algoritma, pseudocode, atau lainnya), bagaimana menguji dan mengumpulkan data. Deskripsi alur penelitian harus didukung oleh referensi sehingga penjelasan yang ditulis dapat diterima secara ilmiah.

Tabel dan gambar berikut dibuat rata tengah (center) seperti yang ditunjukkan pada tabel dan gambar di bawah dan dirujuk pada naskah artikel.

\begin{table}[h]
\centering
\caption{Performa dari mesin X.}
\begin{tabular}{lcc}
Variable & Speed (rpm) & Power (kW) \\
x & 10 & 8.6 \\
y & 15 & 12.4 \\
z & 20 & 15.3 \\
\end{tabular}
\label{tab:1}
\end{table}

Nama tabel berada di atas tabel dengan menggunakan font Times New Roman ukuran 10pt seperti yang ditampilkan pada Tabel \ref{tab:1}. Nama tabel menggunakan spacing after atau spasi setelah baris teks sebesar 6pt. Tabel hanya menggunakan border horizontal dengan ukuran font 8pt.

\begin{figure}[h]
\centering
\framebox[0.8\textwidth][c]{Gambar 1. Efek memilih pengalihan berbeda pada kondisi dinamis.}
\label{fig:1}
\end{figure}

Nama gambar berada di bawah gambar dengan menggunakan font Times New Roman ukuran 10pt seperti yang ditampilkan pada Gambar \ref{fig:1}. Nama gambar menggunakan spacing before atau spasi sebelum baris teks sebesar 6pt.

Persamaan dibuat menggunakan rata tengah dengan penomoran yang menggunakan rata kanan pada satu baris. Penomoran menggunakan tanda kurung ( ). Ketika disebutkan pada artikel, penyebutan dapat dilakukan dengan cara menyebutkan Persamaan \ref{eq:1}.
\begin{equation}
c2 = a2 + b2 \label{eq:1}
\end{equation}

\section{Hasil dan Pembahasan}
Bagian ini membahas hasil dari penelitian dan pada waktu yang sama juga memberikan pembahasan dan yang komprehensif. Hasil penelitian dapat disajikan menggunakan gambar, grafik, tabel, dan lainnya yang membuat pembaca dapat memahami hasil penelitian dengan mudah. Pembahasan dapat dibuat dengan menggunakan beberapa sub-bab.

\subsection{Sub Bab 1}
Lorem ipsum dolor sit amet, consectetur adipisicing elit, sed do eiusmod tempor incididunt ut labore et dolore magna aliqua.

\subsection{Sub Bab 2}
Lorem ipsum dolor sit amet, consectetur adipisicing elit, sed do eiusmod tempor incididunt ut labore et dolore magna aliqua.

\subsection{Sub Bab 3.2}
Lorem ipsum dolor sit amet, consectetur adipisicing elit, sed do eiusmod tempor incididunt ut labore et dolore magna aliqua.

\section{Kesimpulan}
Pada bagian ini, penulis memberi pernyataan mulai dari apa yang diharapkan dari penelitian, yang ditulis pada bagian “Pendahuluan”, sampai dengan hasil yang diperoleh pada bagian “Hasil dan Pembahasan”, sehingga menjadi sebuah kesatuan yang dapat dijelaskan secara singkat, padat, dan jelas. Pada bagian ini juga dapat ditambahkan mengenai rencana penelitian berikutnya berdasarkan hasil yang diperoleh.

Berdasarkan hasil penelitian, dapat disimpulkan bahwa:
\begin{enumerate}
\item Adversarial traffic dengan randomized JA3 fingerprints dapat meningkatkan evasion rate terhadap NIDS berbasis signature dan anomaly.
\item Serangan FGSM secara signifikan mengurangi F1-score model XGBoost, CatBoost, dan MLP.
\item Fitur terkait volume paket (\texttt{total\_packets}, \texttt{avg\_packet\_size}) lebih kritis daripada fitur JA3 dalam performa model.
\item Adversarial training meningkatkan robustness model sebesar 20\% pada average.
\end{enumerate}
Rencana penelitian berikutnya termasuk integrasi dengan sistem NIDS nyata seperti Zeek atau Suricata, serta penyelidikan terhadap serangan yang lebih kompleks seperti DeepFool dan CW dalam domain jaringan.

\section{Daftar Pustaka}
Referensi utama berupa jurnal dan proceeding internasional. Semua referensi sebaiknya merupakan referensi yang berkaitan dengan penelitian dan bersumber dari penelitian terbaru, biasanya lima tahun terakhir. Referensi ditulis menggunakan gaya IEEE. Mohon gunakan format yang konsisten dalam penulisan referensi atau daftar pustaka. Mohon gunakan reference manager Mendeley untuk mengelola daftar pustaka dan sitasi.

\begin{thebibliography}{99}
\bibitem{1}
J.~K.~Author, ``Title of chapter in the book,'' in \textit{Title of His Published Book}, xth~ed. City of Publisher, Country if not USA: Abbrev. of Publisher, year, ch.~x, sec.~x, pp.~xxx--xxx.

\bibitem{2}
L.~Stein, ``Random patterns,'' in \textit{Computers and You}, J.~S.~Brake, Ed. New York: Wiley, 1994, pp.~55-70.

\bibitem{3}
J.~Jones. (1991, May 10). Networks (2nd~ed.) [Online]. Available: http://www.atm.com

\bibitem{4}
D.~Casadei, G.~Serra, and K.~Tani, ``Implementation of a Direct Control Algorithm for Induction Motors Based on Discrete Space Vector Modulation,'' \textit{IEEE Transactions on Power Electronics}, vol.~15, no.~4, pp.~769--777, 2007.

\bibitem{5}
R.~W.~Sperry, ``Science, values, and survival,'' \textit{Journal of Humanistic Psychology}, vol.~26, no.~2, pp.~8-24, Spring 1986. doi:10.1177/0022167886262002
\end{thebibliography}
\end{document}